\documentclass[12pt]{article}
\usepackage{notes-project}
\usepackage{graphicx}
\usepackage{subcaption}
\usepackage{listings}

\graphicspath{{lec04/}}
\numberwithin{figure}{section}

\lstset{
  language=Python,
  basicstyle=\ttfamily\footnotesize,
  frame=single,
  framesep=6pt,
  xleftmargin=6pt,
  xrightmargin=6pt,
  breaklines=true,
  breakatwhitespace=true,
  showstringspaces=false,
  columns=fullflexible,
  keepspaces=true,
  keywordstyle=\color{blue!60!black},
  commentstyle=\color{gray},
  stringstyle=\color{green!40!black},
}

\begin{document}

\begin{center}
  {\LARGE\bfseries Lecture 4: Random Process}\\[0.25em]
  {\large Introduction to Data Analysis Techniques}\\[0.15em]
\end{center}

\section{Random Variable} 

\section{Probability Density Function}

\section{Random Number Generator} 

\begin{dfn}[Pseudo Random Number Sequence]
A \emph{pseudo random number sequence} is a sequence which looks random 
\end{dfn} 

\begin{dfn}[Pseudo Random Number Generator (PRNG)]
A \emph{pseudo random number generator} is a deterministic algorithm which generates a pseudo random number sequence. 
\end{dfn} 

\begin{example}[Simple Generator]
\[x_n = (a x_0 + c) \% m\]
\end{example}


\end{document}
